% SEGMENT OLP-0471-S01
% Part: normal-modal-logic
% Chapter: sequent-calculus
% Section: introduction

\documentclass[../../../include/open-logic-section]{subfiles}

\begin{document}

\olfileid{nml}{seq}{int}

\olsection{அறிமுகம்}

கூற்றுத் தருக்கத்திற்கான தொடரணிக் கணியத்தைக்
\iftag{prvBox}{$\Box$\iftag{prvDiamond}{ மற்றும்~}{}}{}%
\iftag{prvDiamond}{$\Diamond$}{} ஆகியவற்றைக் கையாளும் கூடுதல் விதிகளால்
விரிவாக்கலாம். எடுத்துக்காட்டாக, $\Log{K}$ க்கு \Log{LK} உடன்
பின்வருவனவும் உள்ளன:
\[
  \iftag{prvBox}{\iftag{prvDiamond}{% <> and [] primitive
    \Axiom$\Gamma \fCenter \Delta, !A$
    \RightLabel{$\Box$}
    \UnaryInf$\Box\Gamma \fCenter \Diamond\Delta, \Box!A$
    \DisplayProof
    \qquad
    \Axiom$!A, \Gamma \fCenter \Delta$
    \RightLabel{$\Diamond$}
    \UnaryInf$\Diamond!A, \Box\Gamma \fCenter \Diamond\Delta$
    \DisplayProof
}{% only Box primitive
\Axiom$\Gamma \fCenter !A$
\RightLabel{$\Box$}
\UnaryInf$\Box\Gamma \fCenter \Box!A$
\DisplayProof
}}{% only <> primitive
  \Axiom$!A \fCenter \Delta$
  \RightLabel{$\Diamond$}
  \UnaryInf$\Diamond!A \fCenter \Diamond\Delta$
  \DisplayProof
}
\]
$\Log{K}$ இன் விரிவாக்கங்களுக்கு மேலும் கூடுதல் விதிகளையும் சேர்க்க
வேண்டும்.

ஒவ்வொரு வாய்ப்புநிலைத் தருக்கத்திற்கும் இத்தகைய ஒரு தொடரணிக் கணியம்
இருப்பதில்லை. பொருண்மையியல்படி எளியதான $\Log{S5}$ ஐ அணுகுறவுகளையே
பயன்படுத்தாமல் வரையறுக்கலாம்; எனினும், \Cut{} விதி இல்லாமல் நிறைவானதாக
இருக்கும்படி $\Log{LK}$ இலிருந்து கிடைக்கும் ஒரு தொடரணிக் கணியம் அதற்கு
இருப்பதாக அறியப்படவில்லை. ஆனால் அதற்கு வெட்டற்ற நிறைவான
\emph{மிகைத்தொடரணிக்} கணியம் ஒன்று உள்ளது.

\end{document}
